\begin{frame}
	\frametitle{Hardware e limitações}
	\begin{columns}
		\column{.5\linewidth}
		\begin{beamercolorbox}[sep=5pt,rounded=true,shadow=true]{title}
			A promessa de hardware
		\end{beamercolorbox}
		\begin{itemize}
			\item Sem pesos em ponto flutuante, não há necessidade de multiplicadores de ponto flutuante --- só somadores de inteiros.
			\item Abre espaço para ASICs dedicados, mais simples e mais eficientes por operação.
		\end{itemize}
		\column{.5\linewidth}
		\begin{beamercolorbox}[sep=5pt,rounded=true,shadow=true]{title}
			Os limites atuais
		\end{beamercolorbox}
		\begin{itemize}
			\item Kernels/\textit{runtimes} para pesos ternários ainda são imaturos em GPUs convencionais, otimizadas para INT8/FP16 denso.
			\item Exige treinar do zero: não há como converter um checkpoint FP16 já treinado com a mesma qualidade.
			\item Ganho de paridade só aparece a partir de $\sim$3B parâmetros.
		\end{itemize}
	\end{columns}
\end{frame}
